\section{Part VIII --- What this prototype does not (yet) do}

Be the first to say these things; a committee respects a candidate who knows the limits of the code.

\begin{enumerate}[leftmargin=*]
\item \textbf{Batching efficiency.} PhyloGFN updates the features of a whole batch with one batched \code{einsum}; the network environment loops over the trajectories of a batch in Python (\code{batch\_apply\_actions}) because their cases differ. On the 6-taxon test the network sampler is about $2.5\times$ slower per epoch than the tree sampler for the same batch. The fix is mechanical (group rows by \code{case}, batch the four cases) and does not change any result.
\item \textbf{Only the continuous branch-length variant} (PhyloGFN-Continuous, App.~F) is ported --- it is the state-of-the-art variant and the repository's default. The discrete-bin edge model, the parsimony mode, DDP multi-GPU training and AMP are not ported.
\item \textbf{Rooted terminal objects.} PhyloGFN's terminal object is an unrooted tree (last-step parent count $2N-3$, prior $1/(2N-5)!!$). A network is directed, so we keep the root; consequently the tree-mode MLL uses the uniform prior over \emph{rooted} trees. Under a reversible model both priors give the same evidence (each unrooted tree has $2N-3$ rooted versions of equal likelihood), so the two estimators target the same number.
\item \textbf{$\lambda_R$ is a fixed hyper-parameter.} The proposal (training-and-prior slide) conditions the policy on $\lambda_R$ and places a hyper-prior on it; here $\lambda_R$ is a config value and the posterior over $R$ is conditional on it. Two runs with different $\lambda_R$ give different posteriors; this is expected and is the reason the hyper-prior is scheduled. Moreover the proposal's $p(G)\propto e^{-\lambda_RR}$ \emph{over the class} carries a hidden preference for large $R$ through the class sizes $r(N,R)$ (Part~IV.6); the \code{UNIFORM\_R} option makes the penalty per reticulation explicit and is the version to recommend to the supervisors (Q11 of the revision plan).
\item \textbf{Per-site mixture only.} The reward implements the per-site displayed-tree model (Jin et al.\ 2006). The per-locus variant (segmental transfer: whole blocks follow one displayed tree) needs the mixing sum inside the product over loci instead of over sites --- a change of one loop in \code{\_merge\_features} (mix per block, not per column) plus a block map, scheduled for the recombination data.
\item \textbf{No incomplete lineage sorting.} As in the proposal: the displayed-tree likelihood does not model ILS; ILS-only data are a model-confounder measurement, not a sampler target.
\item \textbf{$\theta$ is drawn at the \textsc{Reticulate} step}, i.e.\ Phase 1 and Phase 2 of the proposal are interleaved exactly as PhyloGFN interleaves branch lengths with merges. The terminal object and the reward are the same; the state graph differs only in \emph{when} continuous values are attached. Both are valid GFlowNet state spaces; the interleaved one reuses PhyloGFN's machinery unchanged.
\item \textbf{Trajectory balance only.} Subtrajectory balance and the learned backward policy of the proposal are not in this prototype; the parent-count function (\code{num\_parents}) is exactly what a learned $\PB$ would have to normalise over, so the hook is in place.
\item \textbf{No $k$-NN candidate filtering.} The policy still scores all $\binom{l}{2}$ pairs plus $l$ singles; the filtering (objective O2) is a separate contribution and is orthogonal to the network extension.
\item \textbf{Small-scale evidence only.} Everything here was verified on 3--8 taxa (exhaustively) and trained for minutes on a CPU on 6-taxon simulations. The DS1--DS8 runs with the full PhyloGFN architecture are the next step and need a GPU.
\end{enumerate}
\clearpage
